\chapter{Rotation II: angular momentum and rolling}\label{ch:rot2}

Linear momentum had its chapter; angular momentum is its equal partner ---
the second conserved quantity of mechanics and the reason gyroscopes,
planets, and skaters behave the way they do. The bookkeeping mirrors
chapter~3 exactly, with one new subtlety: the origin you measure from
matters.

\section{Angular momentum: definition}\label{sec:ang-l}

For a particle,
\begin{equation}
  \mathbf L=\mathbf r\times\mathbf p\label{eq:ldef}
\end{equation}
--- momentum filtered through the position relative to a chosen origin
(figure~\ref{fig:lvec}). A straight-line mover carries nonzero $\mathbf L$
about any origin off its line --- the quantity is about the
\emph{rotation of the line of sight}, not about actual circling.

\begin{figure}[htbp]\centering
\begin{tikzpicture}[scale=1.1,force/.style={->,thick,accent}]
\fill (0,0) circle (0.05) node[below left,font=\scriptsize]{origin};
\draw[force] (0,0) -- (3,1.2) node[midway,below,font=\scriptsize]{$\mathbf r$};
\fill (3,1.2) circle (0.05);
\draw[force] (3,1.2) -- (4.5,0.9) node[right,font=\scriptsize]{$\mathbf p$};
\draw[dashed] (3,1.2) -- (4.3,1.72);
\draw[dashed,amber] (0,0) -- (2.6,1.04);
\draw[->,very thick,amber] (1.5,0.6) -- (1.5,0.1) node[below,font=\scriptsize]{$\mathbf L=\mathbf r\times\mathbf p$};
\node[font=\scriptsize] at (4.8,1.8) {out of page};
\end{tikzpicture}

\caption{$\mathbf L=\mathbf r\times\mathbf p$: magnitude $rp\sin\theta$,
direction by the right-hand rule (out of the page here). Notice it needs an
origin --- the same $\mathbf p$ has different $\mathbf L$ about different
points.}
\label{fig:lvec}
\end{figure}

\section{Rigid bodies: when $L=I\omega$ holds}\label{sec:ang-rigid}

Summing \eqref{eq:ldef} over a rigid body rotating about a fixed
symmetry axis gives
\[
  L=I\omega.
\]
The fine print: valid about the fixed axis \emph{or} about the CM --- the
two origins where the cross terms vanish. A rolling wheel's $\mathbf L$ is
$I\omega$ about its CM \emph{plus} the orbital $M(\mathbf R\times\mathbf V)$
of the CM --- both exist, both count.

\section{The rotational second law}\label{sec:ang-taul}

\begin{equation}
  \boldsymbol\tau_{\rm ext}=\frac{d\mathbf L}{dt}\label{eq:taul}
\end{equation}
about an inertial fixed point or about the CM --- the two origins where the
derivation works. Choosing an accelerating origin is the classic trap: the
equation picks up fictitious-torque terms that ruin it.

\section{Conservation of angular momentum}\label{sec:ang-cons}

\begin{principle}[Angular momentum conservation]
$\mathbf L$ is conserved when the net \emph{external torque} is zero ---
weaker than zero external force. A planet feels a huge central force that
exerts no torque, so its $\mathbf L$ (and Kepler's second law) survives.
\end{principle}

\begin{example}[The skater]\label{ex:skater}
A skater spins at $\omega_0=2$\,rad/s with $I_0=3$\,kg\,m$^{2}$, then pulls
her arms to $I=1$\,kg\,m$^{2}$. Find the new spin rate and the energy
change.
\par\smallskip
\textbf{Strategy.} No external torque about the spin axis: $L$ conserved.
$K$ is not --- her arms do work.
\begin{solution}
$\omega=I_0\omega_0/I=6$\,rad/s. Energies:
$K_0=\tfrac12(3)(4)=6$\,J, $K=\tfrac12(1)(36)=18$\,J --- $K$ tripled. The
extra $12$\,J came from her muscles pulling mass inward against the
centrifugal pull: the audit of ch4 applies here too. \checkmark
\end{solution}
\end{example}

\begin{figure}[htbp]\centering
\begin{tikzpicture}[scale=1.0,font=\scriptsize]
\begin{scope}
\draw[fill=gray!20] (0,0) ellipse (0.25 and 0.7);
\draw[thick] (-0.9,0.3) -- (0.9,0.3);
\draw[thick] (-0.9,0.5) -- (0.9,0.5);
\fill[fill=amber!40] (-1.0,0.2) rectangle (-0.8,0.6);
\fill[fill=amber!40] (0.8,0.2) rectangle (1.0,0.6);
\node at (0,-1.0) {arms out: $I$ big, $\omega$ small};
\draw[->,dashed] (0,-0.7) arc (-90:200:0.4);
\end{scope}
\begin{scope}[xshift=3.6cm]
\draw[fill=gray!20] (0,0) ellipse (0.2 and 0.7);
\draw[thick] (-0.3,0.4) -- (0.3,0.4);
\fill[fill=amber!40] (-0.35,0.25) rectangle (-0.15,0.55);
\fill[fill=amber!40] (0.15,0.25) rectangle (0.35,0.55);
\node at (0,-1.0) {arms in: $I$ small, $\omega$ big};
\draw[->,dashed] (0,-0.7) arc (-90:290:0.3);
\end{scope}
\node at (1.8,-1.9) {$L=I\omega$ fixed; $K=\tfrac12I\omega^{2}$ grows --- the arms' work pays it};
\end{tikzpicture}

\caption{The skater trade: $L$ fixed, $K$ tripled --- arms supply the
difference.}
\label{fig:skater}
\end{figure}

\begin{example}[Bullet into a pivoted rod]\label{ex:bulletrod}
A $10$\,g bullet at $400$\,m/s embeds in the tip of a $1$\,kg,
$1$\,m rod pivoted at its top ($I=\tfrac13ML^{2}$). Find the angular
velocity just after.
\par\smallskip
\textbf{Strategy.} Linear momentum is \emph{not} conserved (the pivot
pushes); angular momentum about the pivot \emph{is} (the pivot force has
no lever arm). Use the collision instant only.
\begin{solution}
$L_i=mvL=0.01\times400\times1=4$\,kg\,m$^{2}$/s.
$I_{\rm tot}=\tfrac13(1)(1)+0.01(1)^{2}=0.343$\,kg\,m$^{2}$.
$\omega=4/0.343=11.7$\,rad/s. $K$ audit: before $800$\,J, after
$\tfrac12I\omega^{2}=23.4$\,J --- $97\%$ lost to the embedding, as
inelastic as collisions get. \checkmark
\end{solution}
\end{example}

\begin{figure}[htbp]\centering
\begin{tikzpicture}[scale=1.0,force/.style={->,thick,accent}]
% bullet embedding in rod
\draw[very thick] (0,0) -- (0,3.2);
\fill (0,2.4) circle (0.05) node[right,font=\scriptsize]{pivot};
\draw[force] (-2.2,2.9) -- (-0.2,2.9) node[left,midway,above,font=\scriptsize]{bullet $\mathbf p$};
\fill[fill=gray!30] (-0.15,2.7) rectangle (0.15,3.1);
\draw[->,dashed] (0.5,2.9) arc (0:120:0.7);
\node[font=\scriptsize] at (1.3,3.4) {$\omega$ after};
\node[font=\scriptsize,align=left] at (3.4,1.4) {$L_{\rm in}=p\,d$ at impact\\$p$ NOT conserved:\\pivot pushes};
\end{tikzpicture}

\caption{Bullet-and-rod: conserve $L$ about the pivot, not $p$ --- the
pivot's impulse is external to translation but torque-free about itself.}
\label{fig:rxp}
\end{figure}

\begin{exercise}[easy]
In example~\ref{ex:bulletrod}, which quantities survive: total $\mathbf L$
about the pivot, total $\mathbf p$, total $K$? One-line each. (AP-C)
\end{exercise}

\section{Rolling dynamics, three ways}\label{sec:ang-roll}

A roller on an incline can be solved by forces and torque
($ma=mg\sin\theta-f$, $I\alpha=fr$), by energy
($v^{2}=2gh/(1+\beta)$ for $I=\beta mr^{2}$), or by the instantaneous-axis
trick (torque about the contact point only). All three agree; the energy
route is the fastest.

\begin{example}[The incline race]\label{ex:race}
Hoop ($\beta=1$), disk ($\tfrac12$) and sphere ($\tfrac25$) race down a
$2$\,m drop. Order and margins?
\par\smallskip
\textbf{Strategy.} $v=\sqrt{2gh/(1+\beta)}$ --- smaller $\beta$ finishes
first, independent of mass and radius.
\begin{solution}
$v_{\rm sph}=5.29$, $v_{\rm disk}=5.11$, $v_{\rm hoop}=4.43$\,m/s ---
the code listing below prints the same numbers. The hoop loses because
\emph{all} its mass sits at the rim: half its energy is rotation at every
instant (figure~\ref{fig:race}). \checkmark
\end{solution}
\end{example}

\begin{figure}[htbp]\centering
\begin{tikzpicture}[scale=0.8,font=\scriptsize]
\foreach \x/\lab/\tr/\rt in {0/slide/4/0, 2.4/sphere/2.85/1.15, 4.8/disk/2.65/1.35, 7.2/hoop/2/2}{
  \begin{scope}[xshift=\x cm]
    \draw[thick] (0,0) rectangle (1.8,2.8);
    \fill[accent!70] (0,0) rectangle (0.8,\tr*0.6);
    \fill[amber!60] (1.0,0) rectangle (1.8,\rt*0.6);
    \node at (0.9,-0.4) {\lab};
    \node[rotate=90] at (0.4,\tr*0.3) {$K_t$};
    \node[rotate=90] at (1.4,\rt*0.3) {$K_r$};
  \end{scope}}
\node at (3.6,3.5) {same $mgh$: rotation share slows the finish --- hoop pays half};
\end{tikzpicture}

\caption{Energy partition in the race: same $mgh$, different splits ---
the hoop spends half on spinning.}
\label{fig:race}
\end{figure}

\begin{lstlisting}[language=Python,caption={\texttt{code/ch06\_rolling\_race.py} --- the incline race for several shape factors.},label={lst:race}]
v = math.sqrt(2*G*H/(1+beta))
\end{lstlisting}

Output of a real run:
\begin{lstlisting}[numbers=none]
incline drop h=2.0 m
shape                   beta   v (m/s)    K_tr   K_rot
point mass (slide)       0.0     6.261   19.60    0.00
solid sphere             0.4     5.292   14.00    5.60
solid disk               0.5     5.112   13.07    6.53
hoop                     1.0     4.427    9.80    9.80

finish order: point mass (slide) > solid sphere > solid disk > hoop
sphere beats hoop by 19.5% of hoop's speed
\end{lstlisting}

Notice the partition column: the hoop's $K_{\rm rot}$ equals its
$K_{\rm tr}$ --- $\beta=1$ means rotation and translation split the
budget evenly, and mass/radius never appear. Figure~\ref{fig:speedbeta}
plots the finish speeds against $\beta$: the points fall exactly on
$v=\sqrt{2gh/(1+\beta)}$, so the whole race reduces to one parameter.

\begin{figure}[htbp]\centering
\begin{tikzpicture}
\begin{axis}[width=8.8cm,height=5.2cm,xlabel={shape factor $\beta$ of $I=\beta mr^{2}$},
ylabel={finish speed (m/s)},
xmin=-0.08,xmax=1.08,ymin=0,ymax=7,font=\small,legend pos=north east,
xtick={0,0.4,0.5,1.0},xticklabels={slide,sphere,disk,hoop}]
\addplot[very thick,accent,mark=*,mark size=2.2pt] coordinates
{(0,6.261) (0.4,5.292) (0.5,5.112) (1.0,4.427)};
\addlegendentry{race simulation}
\addplot[dashed,domain=-0.08:1.08,samples=2] {6.261/sqrt(1+x)};
\node[font=\scriptsize,anchor=west] at (axis cs:0.05,3.6) {$v=\sqrt{2gh/(1+\beta)}$};
\node[font=\scriptsize,accent,anchor=west] at (axis cs:0.42,6.0) {$K_{\rm rot}$ tax grows with $\beta$};
\end{axis}
\end{tikzpicture}

\caption{Finish speed versus the shape factor $\beta$: simulated points
on the analytic curve. Bigger $\beta$ means a bigger rotational tax ---
and nothing else matters.}
\label{fig:speedbeta}
\end{figure}

\begin{example}[Two-disk angular collision]\label{ex:twodisk}
A disk ($I_1=2$\,kg\,m$^{2}$) spinning at $10$\,rad/s drops onto a
stationary disk ($I_2=3$\,kg\,m$^{2}$) on the same axle; friction locks
them. Find the common $\omega$ and the energy lost.
\par\smallskip
\textbf{Strategy.} Angular-momentum conservation about the axle --- the
angular version of the perfectly inelastic collision.
\begin{solution}
$\omega=I_1\omega_0/(I_1+I_2)=20/5=4$\,rad/s.
$K_i=\tfrac12(2)(100)=100$\,J;
$K_f=\tfrac12(5)(16)=40$\,J --- $60\%$ lost, the angular mirror of
chapter 3's inelastic hit. \checkmark
\end{solution}
\end{example}

\begin{example}[Child on a turntable]\label{ex:turntable}
A $40$\,kg child stands at the rim of a $200$\,kg, $3$\,m-radius
turntable ($I=\tfrac12MR^{2}=900$\,kg\,m$^{2}$) spinning at
$\omega_0=0.5$\,rad/s, then walks to the centre. Find the new $\omega$.
\par\smallskip
\textbf{Strategy.} No external torque about the axis: $I\omega$ conserved,
with the child as a point mass whose $I$ shrinks.
\begin{solution}
$I_0=900+40(9)=1260$\,kg\,m$^{2}$; at centre $I_1=900+0$.
$\omega_1=0.5\times1260/900=0.70$\,rad/s --- the same skater effect with
the child playing the arms. $K$ rises from $158$ to $220$\,J, paid by the
child's steps. \checkmark
\end{solution}
\end{example}

\begin{exercise}[medium]
A $40$\,kg child walks from rim to centre on a $200$\,kg, $3$\,m-radius
turntable ($I=\tfrac12MR^{2}$) spinning at $0.5$\,rad/s. Find the new
$\omega$. (AP-C FR-style)
\end{exercise}

\section{Gyroscopic precession: the order of magnitude}\label{sec:ang-gyro}

A spinning top tilted at angle $\theta$ has $\mathbf L$ along its axis;
gravity's torque $\tau=mgd\sin\theta$ is perpendicular, so $\mathbf L$
\emph{rotates} rather than tips (figure~\ref{fig:gyro}). The precession
rate follows from $\Delta L=\tau\Delta t$ geometry:
\[
  \Omega=\frac{\tau}{L}=\frac{mgd}{I\omega}
\]
--- slower spin, faster wobble. That's why tops precess more as they wind
down.

\begin{figure}[htbp]\centering
\begin{tikzpicture}[scale=1.0,force/.style={->,thick,accent}]
\draw[thick] (0,0) -- (2.5,0);
\fill (0,0) circle (0.05) node[below,font=\scriptsize]{pivot};
\draw[thick] (2.5,0) ellipse (0.25 and 0.7);
\draw[force] (2.5,0) -- (2.5,-0.9) node[below,font=\scriptsize]{$mg$};
\draw[force,very thick,amber] (1.2,0) -- (3.2,0) node[right,font=\scriptsize]{$\mathbf L$};
\draw[force,amber] (0,0.3) -- (0,0.9) node[above,font=\scriptsize]{$\boldsymbol\tau$};
\draw[->,dashed] (1.2,-1.6) arc (0:50:1.2) node[midway,right,font=\scriptsize]{precession};
\node[font=\scriptsize,align=left] at (4.6,0.4) {$\Omega=\frac{\tau}{L}$:\\gravity makes\\it precess};
\node[font=\scriptsize,align=left] at (4.6,0.4) {};
\end{tikzpicture}

\caption{Precession bookkeeping: $\boldsymbol\tau$ perpendicular to
$\mathbf L$ turns the vector instead of shrinking it. $\Omega=\tau/L$ sets
the rate.}
\label{fig:gyro}
\end{figure}

\section{In engineering practice: reaction wheels and flywheels}\label{sec:ang-eng}

Satellites steer with \emph{reaction wheels} --- small spun disks whose
speed changes trade $\mathbf L$ with the bus, no fuel spent: this chapter's
conservation law is the whole attitude-control budget. Flywheel energy
storage is the same $K=\tfrac12I\omega^{2}$ as grid batteries: a
$500$\,kg carbon-fiber wheel at $20{,}000$\,rpm banks megajoules in
rotation, which is exactly the hoop's tax of
figure~\ref{fig:race} run to its useful extreme --- mass at the rim is
where you want it.

\section{Common mistakes}\label{sec:ang-mistakes}

\begin{tripwire}[$L=I\omega$ about the wrong axis]
Valid about the fixed axis or the CM only. A rolling wheel about an
arbitrary point needs $I\omega+\mathbf R\times M\mathbf V$.
\end{tripwire}
\begin{tripwire}[Conserving $L$ through external torque]
The bullet-in-rod survives precisely because the pivot torque about the
pivot is zero --- any other origin, or any other external torque, breaks it.
\end{tripwire}
\begin{tripwire}[Forgetting translational $K$ in rolling]
A rolling object carries $\tfrac12mv^{2}+\tfrac12I\omega^{2}$ --- drop the
first term and every race comes out tied.
\end{tripwire}
\begin{tripwire}[$\tau=dL/dt$ about a non-inertial origin]
The law holds at inertial points and the CM. An accelerating origin
invents torques that aren't there.
\end{tripwire}
\begin{tripwire}[Cross-product sign]
$\mathbf r\times\mathbf p\neq\mathbf p\times\mathbf r$ --- the order flips
the sign. Fix the convention once per problem.
\end{tripwire}

\section{Summary}\label{sec:ang-summary}

\begin{table}[ht]\centering
\begin{tabularx}{\linewidth}{@{}lY@{}}
\toprule
\tablehead\thead{Idea} & \thead{One-line version}\\
\midrule
definition & $\mathbf L=\mathbf r\times\mathbf p$ --- origin-dependent \\
rigid body & $L=I\omega$ about fixed axis or CM \\
dynamics & $\boldsymbol\tau=d\mathbf L/dt$ at inertial origin or CM \\
conservation & survives iff external torque $=0$ --- central forces OK \\
rolling & $v^{2}=2gh/(1+\beta)$; smaller $\beta$ wins \\
skater & $L$ fixed, $K$ grows --- muscles pay \\
gyroscope & $\Omega=\tau/L$; perpendicular torque turns, doesn't tip \\
\bottomrule
\end{tabularx}
\caption{Angular momentum cheat sheet.}
\label{tab:angsummary}
\end{table}

\begin{wins}[Definition of done]
You can now: compute $\mathbf L$ for particles and rigid bodies; state
when conservation applies; solve skater, bullet-rod, and two-disk
collisions; run the incline race three ways; and size gyroscopic
precession to order of magnitude.
\end{wins}

\chapternotes{%
The treatment follows MIT 8.01SC weeks 11--12 plus its gyroscope module
\cite{ocw801} and AP Physics C unit 6 \cite{apced}. Morin chs.~8--9
\cite{morin} is the source of the spool and sliding-to-rolling exercises;
OpenStax ch.~11 \cite{openstax1} covers the same material. The race
numbers are the listing's real output.}

\begin{rulebox}[{Watch instead}]
\begin{enumerate}
\item \emph{Angular Momentum, Torques, Conservation of Angular Momentum}
--- Walter Lewin, MIT 8.01x Lecture 20, 51:15,
\url{https://youtu.be/sNaaL19opxw} --- the whole chapter at full length
(accessed 2026-10-11).
\item \emph{Conservation of angular momentum} --- Khan Academy, 7:16,
\url{https://youtu.be/vfmLI150g4w} --- \S\ref{sec:ang-cons} worked quickly
(accessed 2026-10-11).
\end{enumerate}
\end{rulebox}

\section*{Further reading}
\addcontentsline{toc}{section}{Further reading}

\begin{itemize}
\item MIT OpenCourseWare, \emph{8.01SC Classical Mechanics}, weeks 11--12
--- angular momentum and rolling with recitation solutions. At
\url{https://ocw.mit.edu}, course ``8.01SC'' (verified 2026-10-11).
\item OpenStax, \emph{College Physics 2e}, ch.~10--11 --- free text at
\url{https://openstax.org} (verified 2026-10-11).
\item College Board, \emph{AP Physics C: Mechanics CED}, Unit 6 --- the
exam blueprint for angular momentum. At
\url{https://apstudents.collegeboard.org} (verified 2026-10-11).
\item D.~Morin, \emph{Introduction to Classical Mechanics}, chs.~8--9 ---
angular momentum and rolling at the harder level; ISBN 978-0521876223.
\item The Physics Classroom, \emph{Angular Momentum} tutorials --- readable
remediation, at \url{https://www.physicsclassroom.com} (verified
2026-10-11).
\item Khan Academy, \emph{Moments, torque, and angular momentum} ---
drills matching the whole chapter, at \url{https://www.khanacademy.org}
under Science $\to$ Physics (verified 2026-10-11).
\end{itemize}

\section*{Additional exercises}
\addcontentsline{toc}{section}{Additional exercises}

\begin{exercise}[easy]
Find the magnitude and direction of $\mathbf L$ for a $2$\,kg mass at
$\mathbf r=(3,0,0)$\,m moving at $\mathbf v=(0,4,0)$\,m/s. (AP-C)
\end{exercise}
\begin{exercise}[easy]
Decide whether $L$ is conserved in each: (a) planet orbiting the sun,
(b) figure skater pulling in arms, (c) frisbee spinning while gliding to
a stop, (d) the bullet-rod of example~\ref{ex:bulletrod} \emph{after}
impact. (AP-C)
\end{exercise}
\begin{exercise}[easy]
A rolling ball has $v=5$\,m/s and radius $0.1$\,m. Find $\omega$ and
the contact point's instantaneous speed. (AP-C)
\end{exercise}
\begin{exercise}[easy]
Explain in one sentence why the skater's $K$ increases while $L$ stays
fixed. (F=ma)
\end{exercise}
\begin{exercise}[easy]
A spinning disk has $I=0.5$\,kg\,m$^{2}$, $\omega=20$\,rad/s. Find $L$
and $K_{\rm rot}$. (AP-C)
\end{exercise}
\begin{exercise}[medium]
A $40$\,kg child stands at the rim of a $150$\,kg, $4$\,m-radius
turntable ($I=\tfrac12MR^{2}$) at rest, then walks clockwise at
$1$\,m/s relative to the ground. Which way and how fast does the table
turn? (AP-C FR-style)
\end{exercise}
\begin{exercise}[medium]
Rank these down a $3$\,m incline by finish order and compute each speed:
solid sphere, hollow sphere ($\beta=2/3$), disk, hoop. (F=ma)
\end{exercise}
\begin{exercise}[medium]
Example~\ref{ex:bulletrod} with a bullet hitting \emph{halfway} down the
rod: recompute $\omega$ and the energy loss. (AP-C)
\end{exercise}
\begin{exercise}[medium]
Two angular disks merge: $I_1=1$, $\omega_1=12$ and $I_2=4$ spinning
\emph{oppositely} at $\omega_2=-3$. Find the common $\omega$ and energy
lost. (F=ma)
\end{exercise}
\begin{exercise}[medium]
A $0.5$\,m wrench applies $40$\,N tangentially to a $0.02$\,m-radius
shaft. Find the torque and, for a $0.1$\,kg\,m$^{2}$ flywheel, the
$\omega$ after $5$\,s from rest. (AP-C)
\end{exercise}
\begin{exercise}[medium]
A gyroscope with $I=0.01$\,kg\,m$^{2}$, $\omega=200$\,rad/s, $d=0.1$\,m,
$m=0.5$\,kg: compute the precession rate in rad/s and rpm. (F=ma)
\end{exercise}
\begin{exercise}[hard]
A bowling ball is set sliding with no spin at $v_0$ on a lane with
$\mu_k$. Show it transitions to pure rolling at $v=\frac{5}{7}v_0$ ---
track both $v$ and $\omega$ as friction acts. (USAPhO-style)
\end{exercise}
\begin{exercise}[hard]
A spool of thread lies on a table ($I=\beta mr^{2}$, thread wrapped at
inner radius $r_i<R$). Pulling the thread horizontally: does it roll
toward or away from you? Find the critical angle where the direction
reverses. (APhO-style classic)
\end{exercise}
\begin{exercise}[medium]
Energy split check: for the race of example~\ref{ex:race}, what fraction
of each shape's kinetic energy is rotational? Show $\beta/(1+\beta)$.
(AP-C)
\end{exercise}
\begin{exercise}[medium]
A planet at perihelion ($r_1$, $v_1$) reaches aphelion at $2r_1$. Find
its aphelion speed by $L$ conservation, and explain why this is Kepler's
second law. (F=ma)
\end{exercise}
